\section{2024-12-04}

Spannung mal Elektronenladen = Energie

Also benötigt man jeweils Energie durch $17 V = \Delta U_B$, um diese Energei
in Licht umzuwandeln zu können. Alles Licht hat gleiche Fabe, damit gleiche
Wellenlänge, jedes Photon bekommt also die gleiche Energie $\Delta E = e\cdot
\Delta U_B = 17 eV$.

\startformula
1eV = 1 e \cdot 1 V \approx 1.602 \cdot 10^{-19}J
\stopformula

Für Quecksilber ergäbe sich $\Delta E = 4.9eV$

Atome können also, abhängig von der Sorte, bestimmte Energeimengen aufnehmen. Dann sind sie {\bf angeregt}. Diese Energie wird dann in Form von Licht abgegeben. 
